\section{Experimental Setup}
\label{sec:experiments}

\subsection{Evaluation Protocols}
\label{sec:experiments:protocols}

We use two complementary evaluation pipelines. The first is a frozen
closed-set retrieval protocol, and the second is a separate AnimalCLEF-style
open-set discovery protocol. The two protocols use distinct split
constructions, similarity artifacts, calibration procedures, and evaluation
outputs, and are therefore reported separately. All external test identities
and images are excluded from training and checkpoint selection.

\paragraph{Closed-set retrieval.}
For each external dataset, the gallery is formed from the frozen database
partition and queries are drawn from held-out test observations. Gallery and
query observations are disjoint; images from the same observation are never
split across the two roles. We extract one $\ell_2$-normalized embedding per
image and rank gallery images by cosine similarity. The predicted identity is
the identity of the highest-ranked gallery observation. We repeat the
deterministic gallery/query construction for five prescribed split seeds and
report both per-seed results and a dataset-macro average, so that large
datasets do not dominate the aggregate.

The closed-set split is observation-aware and time-aware. When capture
timestamps or reliable date metadata are available, images from the same
encounter or acquisition sequence are kept together and temporally separated
between gallery and query roles; later observations are used as queries where
the dataset protocol permits this ordering. When timestamps are unavailable,
we use the dataset-specific encounter and filename structure, followed by
visual-similarity grouping for samples without temporal or reliable naming
information. In all cases, a group is assigned wholly to one partition, so
correlated frames, burst images, and near-duplicates cannot be split between
gallery and query. The resulting split manifests and observation-leakage checks are frozen and
reused for every model.

\paragraph{AnimalCLEF-style open-set discovery.}
The AnimalCLEF-style protocol constructs observation-aware roles for each
dataset. It first forms encounter
groups using exact-image matches, reliable dates and encounter tokens, and
dataset-specific naming or folder conventions. Temporal grouping is preferred
when available; only identities lacking temporal or naming structure use
visual-similarity grouping. Approximately 20\% of eligible identities are
assigned to \texttt{query\_novel}, while known identities contribute
encounter-disjoint \texttt{gallery} and \texttt{query\_known} roles. Identities
with only one observation group remain query-only unknowns rather than being
forced into a gallery. The Jaguar Antonio setting is gallery-free and is
evaluated as query-only discovery. WildFusion then produces query--query and
query--gallery similarity matrices using a fixed calibration cohort.
Per-animal clustering configurations are selected on development
similarities, while gallery-free cases use label-free defaults. The selected
configurations are then frozen and applied to held-out similarity matrices
using Chinese Whispers. Clustering thresholds are not refit on evaluation
matrices and the observation-aware splits are not rebuilt.

Thus, the AnimalCLEF-style results measure open-set association and
discovery from similarity matrices, rather than closed-set retrieval from the
model gallery. The final clustering labels are compared with ground truth
only for reporting discovery metrics; they are not used to alter the frozen
split or clustering configurations.

\subsection{Baselines}
\label{sec:experiments:baselines}

We compare four ConvNeXt-Small RGB embedding models using the same frozen
preprocessing, gallery construction, split seeds, and evaluation procedures:
ArcFace (AF), Dynamic ArcFace (DAF), SphereFace2 (SF2), and the proposed
Balanced Hierarchical Category--Animal identity Sub-center SphereFace2
(BHCA-SCSF2). The first three models are loss-function baselines; BHCA-SCSF2
is the proposed hierarchical objective. Unless stated otherwise, each model
is evaluated from its corresponding frozen checkpoint and produces a
2048-dimensional $\ell_2$-normalized embedding. No model is fine-tuned on an
external test dataset.

\subsection{Evaluation Metrics}
\label{sec:experiments:metrics}

For closed-set retrieval, we report Rank-1, Rank-5, Rank-10, CMC@20, CMC@50,
mean average precision (mAP), and identity-balanced Rank-1. Rank-$k$ is the
fraction of queries whose true identity occurs among the first $k$ retrieved
gallery identities. mAP is computed per query over the ranked gallery and
then averaged; identity-balanced Rank-1 first averages within identity and
then across identities. Dataset-macro averages give every external dataset
equal weight.

For open-world evaluation, known-identity retrieval is reported with the same
rank and AP measures on queries whose identities occur in the database.
Novelty detection is evaluated using the frozen knownness score with AUROC,
precision, recall, and F1. End-to-end known recognition additionally requires
both a correct known/novel decision and a correct identity assignment.
Novel-identity association and discovery are evaluated with adjusted Rand
index (ARI), adjusted mutual information (AMI), V-measure, and standard
item-averaged B$^3$ precision, recall, and F1. For online runs we additionally
report reassociation accuracy for repeated observations of discovered
identities. Results are summarized per dataset, openness level, and stream
seed before macro-averaging.

\subsection{Implementation Details}
\label{sec:experiments:implementation}

All models use the same animal-focused RGB preprocessing and their
model-specific frozen checkpoint configuration. Embeddings are normalized
before similarity computation. Closed-set retrieval uses cosine similarity
against the complete frozen gallery. The AnimalCLEF-style protocol uses a
fixed WildFusion calibration cohort and per-animal configurations selected on
development similarities. These configurations are applied to held-out
similarity matrices without refitting thresholds or rebuilding splits. The
closed-set and AnimalCLEF-style pipelines therefore have separate provenance,
calibration, and result artifacts.

The closed-set benchmark uses five split seeds and 1,000 paired bootstrap
replicates for confidence intervals. Paired model comparisons reuse identical
queries and split seeds, and confidence intervals are computed with the
identity-block bootstrap. The AnimalCLEF-style pipeline uses its sealed
observation-aware roles, WildFusion calibration sample, per-animal frozen
configuration, and Chinese Whispers clustering protocol.
We record the checkpoint, preprocessing configuration, split-manifest hashes,
calibrator hashes, random seeds, and software environment with every result
bundle. Numerical values are reported in Sec.~\ref{sec:results} once the
corresponding final-test runs are complete.
